%% ===========================================================================
\section{Analysis and discussion}
\label{sec:discussion}
%% ===========================================================================

\subsection{What the errors tell us}
\label{sec:errors}

The decomposition of \cref{sec:formulation} separates three failure modes that a single F1 score conflates, and the three judges fail in different ways.

\paragraph{Zero-shot Laya lacks the information, not the decision rule} Before fine-tuning, Laya's errors were dominated by false positives and by preferences for $n$-grams that contain the gold value (``on april'' for ``april'', ``flight from charlotte''), type confusions (``april'' proposed for a day number or a city) and occasional wrong abstentions; its probabilities stayed close to 0.5 even when it was right. No threshold or calibration could repair this: on the MultiATIS++ English test set, it answers on 83\% of the unexpressed slots and chooses the right value in only 20\% of the cases where it is present. The model had to learn what ``this slot is expressed in the last message'' means for this task.

\paragraph{LLM judges choose well but do not abstain} On MultiATIS++, the LLM judges choose the right value as often as the fine-tuned judge when it is present, but they fill 15--21\% of the slots that the user did not express. Their errors are mostly over-eager answers: a value of the right type is present in the list, often one that the user did express for \emph{another} slot, and the model selects it despite the explicit ``not specified'' option and the instruction to use it. Without in-domain supervision, the notion of ``expressed in the last message'' is underspecified for a general model, and the explicit option does not suffice to calibrate it. Whether a few in-context examples change this behaviour remains open (\cref{sec:limitations}).

\paragraph{The pipeline's bottleneck is the proposers} With real candidates, recall errors dominate (\cref{sec:rq2}), especially in French, where the generic proposers (lexicon, GLiNER with English-like slot descriptions) are weaker. This is good news for the architecture: errors are attributable, and the cheapest components to improve are the proposers, which can be changed without retraining the judge.

\paragraph{Judge or tagger?} The sequence tagger (\cref{sec:rq2-tagger}) sets the limit of our claims. When values are literal spans of the utterance, as on MultiATIS++ and MTOP, a tagger trained on the same data is more accurate than the pipeline with generic proposers, and cheaper per utterance. The proposer--judge architecture pays off when values are \emph{not} literal spans: references to the dialogue history, values of the context, and canonical values that differ from the words of the user, which are frequent in French Multi3WOZ. The two are complementary rather than competing: the tagger is a strong proposer, and the judge adds what the tagger lacks, namely the choice among sources, abstention and a calibrated confidence for deciding whether to act.

\todo{Manual error analysis of 200 test cases per dataset and language, classified by the authors: boundary errors, inter-slot confusions, implicit values, annotation errors, history references; include a few short examples (MultiATIS++ excerpts within the limits of the LDC licence).}

\subsection{Cost}
\label{sec:cost}

The fine-tuned judge is cheaper in every respect that we can measure. It needs a fraction of one GPU (below 2\,GiB of memory), against one or two full H200 GPUs for the LLMs; it answers a single request 3 to 6 times faster and processes 5 to 38 times more requests per second in batches. Its training is cheap: minutes for single-turn datasets and about one hour per seed for the largest (bilingual, multi-turn) training set on one GPU. The request-level ratio is smaller than we expected (H1 assumed 10): an LLM that generates a single token is fast on an H200, and our Laya implementation was not optimised for latency (no compilation, no padding to fixed lengths). The memory footprint, which lets the judge share a GPU with the proposers, is the more decisive advantage for deployment. Fine-tuning the LLMs themselves on the same items would likely close part of the accuracy gap \citep{aguirre2025finetuning}, but at a training and serving cost that is precisely what the small judge avoids; we leave this comparison out of scope.

\subsection{Why the small judge is better calibrated}
\label{sec:calibration-discussion}

Two features of the fine-tuned judge plausibly explain its better risk--coverage behaviour. First, it is trained with a strictly proper scoring rule on exactly the decision it makes at inference, with ``not specified'' as an explicit, frequently correct option; the log score rewards honest probabilities, including for abstention \citep{gneiting2007scoring}. Second, the decision is read in one forward pass from a head designed for option probabilities, whereas the LLMs' first-token distribution is a by-product of generative training, flattened or sharpened by instruction tuning. Even so, the small judge is over-confident ($T$ up to 4.5) and increasingly so away from its training distribution, so temperature scaling on development data remains necessary. The dominant factor for selective decision is, however, not the scale but the \emph{ranking}: temperature scaling does not change the order of the requests by confidence, and the LLMs' AURC stays high after scaling because their confident answers are often wrong. Recalibration methods designed for LLM judges, such as linear probes on hidden states \citep{radharapu2026calibrating} or prediction-distribution correction \citep{li2025calibraeval}, may narrow this gap and are a natural next comparison.

\subsection{Implications for hybrid systems}
\label{sec:hybrid}

The three-zone decision turns the judge into a component that knows when to defer. With real candidates, the fine-tuned judge decides 39--68\% of the MTOP requests and nearly all MultiATIS++ and Multi3WOZ requests on its own at a 5\% target risk; the remainder can be asked back to the user or sent to a costlier component. A natural design is a cascade in which the small judge handles confident cases and an LLM the uncertain ones \citep{lee2024orchestrallm}. Our results suggest caution: the LLM judges are \emph{less} reliable than the small judge on the hard cases too, and \citet{tikhomirov2026selection} report that an LLM verifier placed behind encoder confidence gates does not improve the cost--quality frontier because of a selection effect. Asking the user may thus be a better use of uncertainty than escalating to an LLM; measuring such a cascade directly is left for future work.

\subsection{Application domains}
\label{sec:applications}

The experiments cover three kinds of applications: air-travel information (MultiATIS++), personal assistants over eleven domains (MTOP: alarm, calling, event, messaging, music, news, people, recipes, reminder, timer, weather) and multi-domain task-oriented dialogue for tourism services (Multi3WOZ: restaurant, hotel, attraction, train, taxi), in English and French. The mechanism we study is, however, not specific to slot filling: a compact judge that selects one of several proposed options or abstains, trained with a proper scoring rule and calibrated so that its confidence can drive a three-zone decision. It plausibly transfers to other frequent, well-defined decisions for which supervision can be collected, which we have not evaluated: selecting the value of a form field or of a structured record among the outputs of extractors (e-mails, tickets, reports); selecting a tool, or the value of a tool argument, in an agent, with an explicit decision to ask for confirmation; and replacing LLM judges that arbitrate between options inside a pipeline, at a fraction of their cost and memory footprint, which also makes on-premise deployment easier. The selective decision matters most where a wrong action costs more than a question to the user, as in health, finance, public administration or bookings. Two conditions learnt from this study apply in every case: in-domain (and in-language) training data are needed, and the proposers must be strong, since the judge cannot recover a value that is not proposed and a dedicated extractor may otherwise be the better choice (\cref{sec:rq2-tagger}).

\subsection{Limitations}
\label{sec:limitations}

\begin{itemize}
  \item \textbf{Single test pass, development-set decisions.} All protocol decisions were taken on development data, and checkpoint selection used the first 300 development examples of each dataset, which also served for those decisions; the test sets were then used once, under a protocol frozen beforehand. Latencies and the descriptive statistics of the candidate lists come from development data.
  \item \textbf{Gold intent and domain.} Requests use the gold intent (or the gold domain on Multi3WOZ), which isolates slot filling from intent detection. The domain filter also leaves out 1.9\% of the Multi3WOZ gold pairs, which caps recall.
  \item \textbf{Scope.} H3 is supported only for the confusion-oriented augmentation, and the sweeps and the \texttt{hard} ablation were run for the Laya judges on MultiATIS++ English only. The latency part of H1 is not supported (ratio 3 to 6 instead of 10). We did not study the two-step decision (H4), unseen-domain splits, renamed or paraphrased schemas, a generative LLM extractor without proposers, or a cascade with deferral to an LLM; and we did not add the tagger's spans to the real candidate lists, although \cref{sec:rq2-tagger} suggests that it is the most promising change.
  \item \textbf{LLM coverage.} Three LLMs from two families were evaluated, zero-shot only, with one prompt; a third family (Mistral-Small-4) could not be served with our vLLM version. Prompt engineering, few-shot prompting or fine-tuning might improve the LLMs.
  \item \textbf{Languages and data.} Only English and French are covered, on public benchmarks that approximate real assistant traffic. MultiATIS++ is narrow (one domain) and cannot be redistributed (see the Data and code availability statement).
  \item \textbf{Bilingual training.} Joint English--French training is unstable with real candidates on Multi3WOZ (wrong abstentions, \cref{sec:rq3}), probably because the checkpoint is selected on English oracle lists; we therefore use monolingual models.
  \item \textbf{Training objective.} We fine-tune with a supervised log score because Laya's original reinforcement-learning training code is not released; the comparison with the original procedure is not possible. With full labels and calibrated probabilities, a threshold on the probability is already the optimal decision rule, so reinforcement learning would add little to the task as posed here. It becomes relevant in three settings that we did not study: (i) learning the act/ignore/ask policy directly under asymmetric costs, for instance by adding the expected reward $\sum_a p(a \mid q)\,R(a, y^{*})$ of the decision to the log score (exact, since the action set is small); (ii) learning from user feedback (confirmations, corrections) instead of gold labels, as an off-policy contextual bandit; and (iii) learning when to ask at the level of a whole dialogue, with a reward that trades task success against the number of questions and a simulated user, e.g.\ on Multi3WOZ. In each case calibration must be monitored, since a reward objective is not a proper scoring rule.
\end{itemize}

\subsection{Threats to validity}
\label{sec:threats}

\paragraph{Internal validity} (i) \emph{Test-set exposure.} The MultiATIS++ English test split was opened several times during early development; results from these runs are labelled exploratory and none is used for the main claims; all later decisions were taken on development data, and the test sets will be used once per frozen configuration. (ii) \emph{Selection criterion.} The checkpoint criterion was changed twice on development data (global loss, then a class-balanced loss, then development F1), because over-confidence distorted the losses: one C3 seed kept an epoch with a lower loss but 663 test false positives against 304--462 for the other seeds. (iii) \emph{Configuration drift.} An early bug (oracle generators not fitted on training data) produced empty or impoverished candidate lists and invalidated the first runs; another series ran on stale code. Both were detected, the affected runs were re-labelled or discarded, and a guard and a code fingerprint in every run now prevent recurrence. (iv) \emph{Latency} was measured on 100 development examples per dataset with batches of size one; it depends on the implementation (inference engine, compilation, padding), which favours the LLMs, served by an optimised engine.

\paragraph{External validity} Public benchmarks are cleaner than real traffic; annotation conventions (span boundaries, nested MTOP slots, Multi3WOZ state differences) shape what counts as correct. Our hand-written slot descriptions were frozen before evaluation, but other descriptions could favour one judge or the other.

\paragraph{Construct validity} With most slots unexpressed, slot F1 mostly measures abstention; we therefore always report choice accuracy, false-positive rate and abstention when the gold value is absent alongside F1, and frame accuracy as a stricter view. Joint goal accuracy for Multi3WOZ is not reported (\cref{sec:metrics}), so our Multi3WOZ numbers are not comparable with dialogue-state-tracking leaderboards.

\paragraph{Contamination} The benchmarks are public and may be part of the LLMs' training data. Contamination would favour the LLMs, so it cannot explain their lower scores; it could however mask an even larger gap. We have no unseen-domain control.

\paragraph{Statistical validity} Training seeds are few (three to five) and LLM judges currently use one seed; intervals reflect test-sample variability, and seed variability is reported separately. The Holm correction is applied within families of comparisons, not across the whole paper.
